\section{Consistency among electronic-state files}
\label{sec-reading}

The electronic states used in a single \progname{} job may be read from one or more external quantum chemistry calculations. When all states are obtained from the same calculation, the single-package form in \script{chap-interfaces}{script-statepack-single} is sufficient. States obtained from separate calculations are supplied through the multiple-package form in \script{chap-interfaces}{script-statepack-multiple}. The separate calculations may in some cases use different electronic-structure methods or even different quantum chemistry programs.

Whether states from different calculations can be used together depends on the operation to be performed on their wave functions. \progname{} therefore checks the consistency of the loaded states according to the requirements of the requested job before the calculation begins.

For operations involving transition density matrices or linear combinations of electronic states, the participating states must be represented by mutually compatible wave functions. In particular, they must have the same wave-function type, number of electrons, AO basis, and configuration space, together with the other basic settings required by the corresponding wave-function representation. These conditions are required because interstate quantities and linear combinations are evaluated directly between the wave functions and therefore require them to share a common electronic-structure representation.

Difference-density calculations require less restrictive consistency. A reference state and a target state may be obtained with different electronic-structure methods and may contain different numbers of electrons. They must nevertheless use the same molecular geometry, atom ordering, and AO basis so that their one-particle density matrices can be represented and compared in the same AO basis. An important example is multiple-excitation FPHD for a charged system, where a separately calculated neutral state may be used as the reference state. The reference and target calculations may then differ in electron number and electronic-structure method, while the geometry, atom ordering, and AO basis must remain consistent. See Section~\ref{sec-fphd} for the reference-state requirements specific to FPHD.

Independent quantum chemistry calculations introduce an additional issue even when their computational settings are otherwise identical. MOs obtained from separate calculations may differ by arbitrary phase factors. After the relevant states have been loaded, \progname{} compares corresponding MOs from different files and determines their relative signs. The excitation amplitudes or configuration coefficients of the associated electronic states are then adjusted consistently with these orbital-phase differences. This places states read from different files under a common orbital phase convention before interstate properties or linear combinations are evaluated.
This phase adjustment does not remove the consistency requirements described above.

In summary, the admissible combination of electronic-state files is determined by the mathematical operation required by the selected job. \progname{} applies the corresponding consistency checks automatically and rejects combinations that do not provide a compatible representation for that operation.